\section{Endpoint Bisection Certificate}
\label{sec:unified-reality-endpoint-bisection-certificate}

The infinite spiral picture is too coarse if it is read as a literal
endpoint of the full Dirichlet polyline.  The zero-relevant object must
be a finite endpoint model equipped with dominance, tail control, and a
distinguished segment on which the zero acts.  The fixed quartic carrier
supplies the finite template: endpoint zeros are produced by active
branch interference after a dominant endpoint window and a Rouch\'e
ledger have been displayed.  For zeta, the corresponding candidate is
the Riemann--Siegel or approximate-functional-equation endpoint window,
where the completed read is split into two finite arms and a controlled
remainder.

\begin{definition}[Normalized zeta segment channel]
\label{def:unified-reality-normalized-zeta-segment-channel}
For
$$
s=\frac12+\delta+i\tau
$$
and $n\ge1$, the \emph{normalized zeta segment channel} is
$$
u_n(s):=n^{1/2-s}=n^{-\delta}e^{-i\tau\log n}.
$$
Its straightened log-length defect is
$$
\ell_n(s):=\log|u_n(s)|=-\delta\log n.
$$
The raw Dirichlet or eta segment may have a baseline length such as
$n^{-1/2}$; the normalized channel removes that baseline and records
only the unitary versus radial component.  Thus the critical line is the
row
$$
|u_n(s)|=1
$$
for every displayed $n$.
\end{definition}

\begin{theorem}[Weighted endpoint energy reads the normal displacement]
\label{thm:unified-reality-weighted-endpoint-energy-reads-normal-displacement}
For $N\ge2$, define
$$
\mathcal E_N^{\zeta}(s)
:=
\sum_{1\le n\le N}
\frac1n
\left(\log|n^{1/2-s}|\right)^2.
$$
If $s=1/2+\delta+i\tau$, then
$$
\mathcal E_N^{\zeta}(s)
=
\delta^2
\sum_{1\le n\le N}\frac{(\log n)^2}{n}.
$$
Consequently, if $\delta\#0$, then
$$
\sup_N\mathcal E_N^{\zeta}(s)=\infty.
$$
\end{theorem}

\begin{proof}
By
\autoref{def:unified-reality-normalized-zeta-segment-channel},
$$
\log|n^{1/2-s}|=-\delta\log n.
$$
Squaring and summing with weight $1/n$ gives the displayed identity.
The positive summands $(\log n)^2/n$ are unbounded in total: for
$N\ge2$,
$$
\sum_{2\le n\le N}\frac{(\log n)^2}{n}
$$
dominates a cofinal subseries obtained by grouping $n$ into dyadic
blocks, and the block contributions diverge.  Multiplication by
$\delta^2>0$ preserves divergence in the located positive readout.
\end{proof}

\begin{definition}[Finite endpoint segment and bisection parameter]
\label{def:unified-reality-finite-endpoint-segment-bisection-parameter}
Let
$$
P_m(s):=\sum_{1\le n\le m}v_n(s)
$$
be a finite partial read in a displayed endpoint model.  The
\emph{endpoint segment} at $m$ is the affine segment
$$
P_{m-1}(s)+\lambda v_m(s),
\qquad 0\le\lambda\le1.
$$
If the segment contains the origin, its \emph{bisection parameter}
$\lambda_m(s)$ is the value satisfying
$$
P_{m-1}(s)+\lambda_m(s)v_m(s)=0.
$$
\end{definition}

\begin{theorem}[Endpoint bisection is endpoint symmetry]
\label{thm:unified-reality-endpoint-bisection-endpoint-symmetry}
Assume that a finite endpoint segment contains the origin.  Then
$$
\lambda_m(s)=\frac12
$$
if and only if
$$
P_m(s)=-P_{m-1}(s).
$$
\end{theorem}

\begin{proof}
The origin condition is
$$
P_{m-1}(s)+\lambda_m(s)v_m(s)=0.
$$
If $\lambda_m(s)=1/2$, then
$$
P_{m-1}(s)=-\frac12v_m(s),
$$
and hence
$$
P_m(s)=P_{m-1}(s)+v_m(s)=\frac12v_m(s)=-P_{m-1}(s).
$$
Conversely, if $P_m(s)=-P_{m-1}(s)$, then
$$
2P_{m-1}(s)+v_m(s)=0,
$$
so the origin is reached at $\lambda=1/2$.
\end{proof}

\begin{definition}[Riemann--Siegel endpoint model row]
\label{def:unified-reality-riemann-siegel-endpoint-model-row}
A \emph{Riemann--Siegel endpoint model row} for a completed zeta read at
height $\tau$ consists of:
$$
N(\tau)\sim\sqrt{\frac{\tau}{2\pi}},
$$
two finite arms of approximate-functional-equation type
$$
\sum_{n\le x}n^{-s},
\qquad
\chi(s)\sum_{n\le y}n^{s-1},
\qquad
xy\sim\frac{\tau}{2\pi},
$$
a symmetric endpoint choice $x\sim y\sim N(\tau)$, and a remainder row.
On the critical line this row is transported to a real finite
Riemann--Siegel read of the form
$$
2\sum_{n\le N(\tau)}
n^{-1/2}\cos(\theta(\tau)-\tau\log n)
$$
plus the displayed remainder.  The row is a candidate zero-relevant
endpoint model only when dominance and tail bounds identify a local
zero island.
\end{definition}

\begin{definition}[Endpoint bisection certificate]
\label{def:unified-reality-zeta-endpoint-bisection-certificate}
For a critical-strip zeta zero $s=1/2+\delta+i\tau$, a
\emph{zeta endpoint bisection certificate} consists of a cofinal family
of canonical finite endpoint models $F_N(s)$ with endpoint segments
$e_N(s)$ and positive endpoint scales $h_N$, together with rows:
$$
\Xi(s)=F_N(s)+T_N(s),
\qquad
\mathsf{RoucheGap}(F_N,T_N),
$$
$$
\mathsf{ZeroRelevantEndpoint}(e_N,s),
\qquad
\lambda_N(s)\to\frac12,
$$
$$
\lambda_N(s)=\frac12
\Longleftrightarrow
|e^{-(s-1/2)h_N}|=1,
$$
and a cofinal growth row such as
$$
h_N\to\infty.
$$
The Rouch\'e row prevents a projected finite chord from being treated as
a zero certificate.  The straightened equivalence says that bisection in
the endpoint segment is the unit-length condition in the normalized
normal coordinate.
\end{definition}

\begin{theorem}[Endpoint bisection certificates are conditionally sufficient]
\label{thm:unified-reality-endpoint-bisection-certificates-conditionally-sufficient}
If every critical-strip zeta zero carries a zeta endpoint bisection
certificate whose limiting bisection row forces
$$
\log|e^{-(s-1/2)h_N}|\to0,
$$
then the RH fixed-half section exists.
\end{theorem}

\begin{proof}
Let $z:\mathcal Z_{\zeta}$ have complex coordinate
$$
s=\frac12+\delta+i\tau.
$$
The certificate supplies positive scales $h_N$ with $h_N\to\infty$ and
the limiting straightened read
$$
\log|e^{-(s-1/2)h_N}|\to0.
$$
But
$$
\log|e^{-(s-1/2)h_N}|=-\delta h_N.
$$
If $\delta\#0$, then the displayed expression cannot converge to $0$
along a cofinal sequence with $h_N\to\infty$.  Hence
$\neg(\delta\#0)$.  Located fixed-half collapse gives $\delta=0$.
Therefore $\Re(s)=1/2$, and
\autoref{thm:unified-reality-rh-fixed-half-section-equivalence} gives
the fixed-half section.
\end{proof}

\begin{corollary}[Endpoint bisection refines local half-split]
\label{cor:unified-reality-endpoint-bisection-refines-local-half-split}
The endpoint bisection certificate refines the local half-split
certificate by specifying a zeta-native finite endpoint model in which
the half-split is to be read.
\end{corollary}

\begin{proof}
\autoref{def:unified-reality-zeta-local-half-split-certificate} requires
a local two-branch zero island, Rouch\'e tail control, length equality,
and a real-defect split formula.  The endpoint bisection certificate
adds a candidate source for that local island: a Riemann--Siegel or
approximate-functional-equation endpoint window with a distinguished
endpoint segment.  Thus the half-split row is no longer attached to an
arbitrary finite chord; it is attached to a completed zeta endpoint
model whose tail and dominance ledgers are part of the certificate.
\end{proof}

\begin{corollary}[The quartic endpoint row is the finite bisection template]
\label{cor:unified-reality-quartic-endpoint-row-finite-bisection-template}
The fixed quartic endpoint-zero row is the finite template for the zeta
endpoint bisection route.
\end{corollary}

\begin{proof}
By
\autoref{thm:unified-reality-quartic-endpoint-zero-branch-interference},
the quartic endpoint zeros are obtained only after the active colliding
branches, endpoint scale, amplitude rows, and Rouch\'e isolation are
displayed.  The resulting shifted cosine zeros are finite half-period
interference points.  The zeta endpoint bisection certificate asks for
the analogous finite endpoint structure in the completed zeta read: not
a bare infinite polyline endpoint, but a zero-relevant endpoint window
with a distinguished bisection segment and a controlled remainder.
\end{proof}

\begin{corollary}[Solenoid fibres do not hide real scale]
\label{cor:unified-reality-endpoint-bisection-solenoid-no-real-scale}
In an endpoint bisection certificate, solenoid-compatible hidden data
may carry phase and prime-register information, but not the real
positive scale $e^{-\delta h_N}$.
\end{corollary}

\begin{proof}
The solenoid route records compact phase together with prime-adic fibre
data, as in
\autoref{def:unified-reality-solenoid-compatible-resonance-lift}.  The
normal endpoint factor
$$
e^{-\delta h_N}
$$
belongs to the noncompact positive real scale direction.  The
straightened equivalence in
\autoref{def:unified-reality-zeta-endpoint-bisection-certificate}
therefore cannot be discharged by placing $\delta$ inside the hidden
solenoid fibre.  The fibre may preserve residue or phase information;
it is not a reservoir for the real radial defect.
\end{proof}
